\scanpage{268}
\mspage{218}
la variante de \textcircled{a} :

\textcircled{c} Pour \unsure{tte} sous-catégorie pleine $\unsure{F|\mathcal F}$ de $M_{/S}$ telle que
\[
\underline{\mathrm{Fib}}_{W}S\subset F\subset M_{/S}
\]
les inclusions induisent des équivalences
\[
\begin{array}{c}
\underline{\rule[-.24ex]{.08ex}{.48ex}\vphantom{F}(\underline{\mathrm{Fib}}_W S)W_S^{-1}\rule[-.24ex]{.08ex}{.48ex}}\\
\Vert\\
\mathrm{Ho}^{\unsure{\mathrm{lc}|\mathrm{li}}}_{W}(S)\\
\text{\deleted{\unsure{moyennant les}}}\\
\text{moyennant a)}
\end{array}
\xrightarrow{\approx} F(W_S)^{-1}
\xrightarrow{\approx}
\begin{array}{c}(M_{/S})(W_{S})^{-1}\\ \Vert\\ \mathrm{Ho}^{!}_{W}(S)\end{array}
\]

Mais je présume que pratiquement, quand on pourra vérifier \textcircled{a}, on pourra vérifier \textcircled{\unsure{b|c}} également.
% slanted
\margin{cf page 240 \unsure{ff|f}}
% slanted
\margin{Il faudrait ici développer un peu plus le formalisme de $\mathrm{Ho}_{W}^{\mathrm{lc}}(S)$ et des fibres homotopiques, sous l'hypothèse \textcircled{a} \textcircled{b} plus haut.}

5. \uline{Axiomatique des $W$-types d'homotopie relatifs non loc. constants} $\mathrm{Ho}^{\unsure{\mathrm{lc}|\mathrm{c}}}_{W}(S)$

On n'a rien dit \unsure{jusqu'en|jusque} ce point d'éventuels \inserted{$W$-}Types d'homotopie relatifs (sur un \unsure{objet} $S$ de $M$) qui ne seraient \uline{pas} ``$\mathrm{loc}^{\mathrm{t}}$ constants''. Il en était question \unsure{pourtant} dans loc. cit (VI, théorème \unsure{salade} 2, partie II du théorème\unsure{)}, pages \unsure{107|109}-113.). Je profite de l'occasion pour revenir là dessus. Il faut introduire, dans chaque $M_{/S}$, un deuxième localiseur
\[
W_S^0\subset\mathrm{Fl}(M_{/S}),\quad W_S^0\subset W_S
\]
(plus fin que $W_S$), ce qui permettra de définir
% archive sheet number: 268
